\section{Error Decomposition and Feature Importance}
\label{sec:errordecomp}

The comparison table of Section~\ref{sec:metrics} reports one number per
model per split. Error analysis (Session 6) decomposes that number along
observable dimensions of the panel: prediction year, country, and
emitter-size tier, so that a claim like ``the model beats persistence''
can be scoped to where it is actually true.

\subsection{Conditional error and within-group skill}

For any partition of the evaluation rows into groups $\mathcal{G}$ (years,
countries, or size tiers), the conditional MAE of model $\mathrm{m}$ in
group $G \in \mathcal{G}$ is

\begin{equation}
\label{eq:group-mae}
\mathrm{MAE}^{\mathrm{m}}_{G} =
\frac{1}{|G|}\sum_{i \in G}\left| y_i - \hat y_i^{\mathrm{m}} \right|,
\end{equation}

and the within-group skill compares each group against persistence's error
in the same group,

\begin{equation}
\label{eq:group-skill}
\mathrm{Skill}_{G} = 1 -
\frac{\mathrm{MAE}^{\mathrm{m}}_{G}}{\mathrm{MAE}^{\mathrm{p}}_{G}},
\end{equation}

so each tier or year is its own honest comparison. Because overall MAE is a
row-share-weighted average of Equation~\ref{eq:group-mae} over groups, a
positive overall skill (Equation~\ref{eq:skill}) is compatible with negative
skill in most groups when the error mass is concentrated in a few of them.
On this project's test years, the headline skill of the best model lives
almost entirely in the giant-emitter tier while every learned model loses
to persistence among small emitters; the overall number alone would have
hidden that scope.

\subsection{Scale dependence of error rankings}

Absolute errors (Equations~\ref{eq:mae}--\ref{eq:medianae}) weight rows by
the size of the miss in Mt; percentage errors (Equation~\ref{eq:mdape})
weight misses by the emitter's own level. The two views can rank the same
countries in reverse: a giant emitter can produce one of the sample's
largest absolute errors at one of its smallest percentage errors, while a
mid-size volatile emitter shows the opposite signature. This is the
scale-dependence problem of forecast-accuracy measures discussed by
\cite{hyndman2006another}, observed here first on the validation years and
re-verified on the held-out test years before being used in any figure.
Percentage errors additionally degenerate as $y_i \to 0$, which is why the
threshold in Equation~\ref{eq:mdape} is part of the metric's definition.

\subsection{Permutation importance and correlated features}

For a frozen fitted model $\hat f$ with evaluation-set error
$L(\hat f)$, the permutation importance of feature $j$ is the increase in
error when column $j$ is randomly permuted across the evaluation rows,
breaking its relationship to the target while preserving its marginal
distribution \cite{breiman2001random}:

\begin{equation}
\label{eq:perm-importance}
\mathrm{PI}_j = \mathbb{E}\!\left[ L\!\left(\hat f;\, X^{(\pi_j)}\right) \right] - L\!\left(\hat f;\, X\right),
\end{equation}

estimated by averaging over several independent permutations $\pi_j$. When
two features carry overlapping information, as this project's co2 level,
lag, and rolling-mean columns do by construction, permuting either column
alone leaves the model able to lean on the other, so both can receive small
values of Equation~\ref{eq:perm-importance} even when their shared signal
is essential. Importances over correlated features are therefore read as
group structure rather than a per-feature ranking; no feature is dropped or
added in response to them here, since the configuration is frozen
post-test.
